\section{Introduction}
Optimizing the elicitation, analysis, and management of software specifications is a foundational objective in software engineering, particularly as system complexity escalates and delivery cycles become increasingly compressed. Software requirements specifications delineate the intended operational capabilities, architectural boundaries, and performance envelopes of software systems. In requirements engineering practice, specifications are broadly bifurcated into functional requirements, which articulate distinct behavioral capabilities and user interactions, and non-functional requirements, which define quality attributes, security policies, reliability constraints, and operational standards [1,2]. The precise categorization of these requirements is critical across downstream development lifecycles, directly informing architectural component partitioning, safety verification, traceability matrix construction, and contractual compliance audits. When non-functional constraints such as latency thresholds, fault tolerance, or encryption protocols are misclassified or overlooked during initial planning phases, software projects frequently encounter architectural rework, severe cost overruns, and systemic system failures.

Despite its recognized importance, manual requirements classification remains labor-intensive, error-prone, and unsustainable for large-scale enterprise repositories that routinely contain thousands of heterogeneous requirement statements. Software specifications are inherently authored in unconstrained natural language, characterized by stylistic heterogeneity, domain-specific terminologies, varying levels of granularity, and extensive semantic ambiguity [3]. Furthermore, software engineers frequently express critical quality expectations through informal descriptions, incomplete clauses, and domain acronyms, obscuring the clear boundaries between functional operations and non-functional quality attributes [4,5]. Consequently, automating requirements classification through natural language processing and machine learning has become an active area of empirical research, aiming to provide robust, reproducible, and scalable categorization mechanisms across diverse industrial projects.

Early computational approaches to requirements classification relied predominantly on classical feature engineering pipelines, extracting lexical tokens, hand-crafted keywords, and term frequency-inverse document frequency vectors for ingestion into traditional supervised classifiers. Although computationally lightweight, these lexical representations remain inherently limited by the high sparsity of software text, the lack of contextual syntactic depth, and an acute vulnerability to vocabulary mismatch [6,7]. With the advent of deep contextual language models and self-attention architectures, researchers transitioned toward fine-tuning bidirectional transformers, such as BERT and its architectural variants, to capture complex syntactic dependencies and contextual nuances within requirements corpora [8]. Sentence-level embedding frameworks, notably Sentence-BERT, further refined semantic vector extraction, yielding fixed-dimensional dense vector manifolds optimized for semantic similarity, semantic search, and thematic clustering.

However, the deployment of large-scale transformer architectures in requirements engineering introduces severe operational constraints. Pre-trained language models demand substantial computational infrastructure for fine-tuning and inference, imposing prohibitive energy and memory overheads that impede routine deployment in lightweight software engineering toolchains [9,10]. Moreover, transformers rely on subword tokenization schemes that frequently fragment domain-specific acronyms, software protocols, and structural keywords, diluting salient discriminative cues. Most critically, software requirements specifications exhibit extreme class imbalance, where common functional requirements vastly outnumber critical quality attributes such as maintainability, portability, or fault tolerance. Deep neural architectures trained on skewed distributions frequently exhibit poor generalization on minority quality classes, experiencing substantial performance degradation when confronted with brief, technical, or colloquially phrased statements.

To overcome the dual challenges of representational sparsity in lexical models and excessive computational complexity in transformer architectures, bio-inspired computing paradigms have emerged as an intriguing alternative for natural language encoding. Natural language and macromolecular biological sequences share fundamental mathematical isomorphisms; both formal human languages and genomic chains function as symbolic sequence systems governed by formal grammars, discrete alphabets, hierarchical dependencies, and non-random informational structures [11]. The informational entropy principles formulated by Shannon reveal that biological nucleotide arrangements encode dense regulatory constraints through character combinations, echoing how syntactic constructs convey operational intent [12]. Furthermore, statistical analyses have confirmed that non-coding DNA sequences exhibit linguistic regularities, such as Zipf's law, demonstrating that genomic arrangements inherently preserve structural and statistical signatures across sequential boundaries [13].

Despite these conceptual parallels, earlier attempts to adapt biological sequence concepts to software requirements classification suffered from severe methodological simplifications. Prior bio-inspired implementations mapped natural language tokens into crude single-nucleotide letters using static, unigram keyword lookup tables. Such static representations treated requirement statements as decoupled character strings, ignoring contextual couplings, syntactic hierarchies, and dependency relationships. By collapsing rich linguistic structures into four isolated characters, these initial frameworks hit an empirical accuracy ceiling of 78.67% on benchmark corpora, demonstrating an inability to resolve fine-grained architectural quality attributes or handle complex requirement phrasing.

To resolve these fundamental limitations, this study presents DNA+, an advanced bio-inspired requirements engineering framework that dynamically translates natural language software requirements into synthetic biological genomic sequences to exploit high-order structural biological grammar for automated classification. Rather than applying static single-nucleotide conversions, DNA+ establishes a dynamic syntactic-semantic dinucleotide codon grammar. In this paradigm, requirement statements are parsed into functional primary anchors, which define the core operational modality, and auxiliary syntactic dependencies, which capture conditional constraints and quality attributes. By mapping these coupled syntactic units onto discrete nucleotide base dimers, the framework preserves structural dependency pairings within a compact biological sequence space.

The synthesized sequences are generated in standard FASTA format, from which an alignment-free genomic physicochemical profile is extracted. Alignment-free oligonucleotide analysis, grounded in classical sequence statistics, evaluates complete sequence compositions without requiring computationally expensive alignment algorithms [14]. The framework extracts monomer frequencies, sixteen dinucleotide relative abundance ratios, GC-content, and Shannon sequence entropy, providing an invariant topological signature for each requirement statement [15]. These structural genomic attributes are subsequently concatenated with classical term frequencies and dense 384-dimensional Sentence-BERT embeddings into a rich 457-dimensional multimodal feature manifold, capturing both abstract contextual semantics and fine-grained structural regularity.

A critical challenge in empirical software engineering classification is the severe imbalance inherent in industrial requirements corpora, where minority quality attributes occur at ratios exceeding 70:1 relative to majority classes [16,17]. Conventional approaches often attempt to alleviate this skew through synthetic minority oversampling techniques, such as SMOTE, which synthesize artificial feature vectors between nearest neighbors. However, empirical investigations have revealed that synthetic oversampling frequently introduces artificial textual artifacts, distorts underlying feature distributions, and impairs the interpretation of software engineering predictors [18,19]. Rather than altering raw empirical distributions, feature space quality and rich multimodal representations dominate algorithmic complexity in mitigating data skew [20]. Consequently, DNA+ is engineered to maintain robust classification boundaries directly on authentic, unmanipulated class distributions.

To establish rigorous experimental validity and avoid common pitfalls identified in empirical software engineering, the validation methodology adheres strictly to established reporting guidelines [21]. Automated predictors evaluated across heterogeneous datasets require standardized non-parametric statistical hypothesis testing and multi-dataset comparison frameworks [22]. To prevent stochastic evaluation bias, cross-validation protocols must incorporate extensive repeated evaluations and non-parametric effect size assessments, mitigating the risk of statistical artifacts [23]. Furthermore, because deep learning and ensemble models are frequently perceived as opaque black boxes, practical adoption requires explainable artificial intelligence techniques that expose internal decision logic. Integrating game-theoretic feature attribution through TreeSHAP provides global interpretability across the multimodal feature space [24], while local model-agnostic perturbations via LIME deliver token-level justification directly understandable by software engineers [25].

The proposed DNA+ framework is benchmarked against contemporary state-of-the-art requirements classification approaches, encompassing modern ensemble stacking architectures [26], comprehensive transformer evaluations [27], adaptive deep learning data augmentation pipelines [28], multilingual back-translation frameworks [29], and domain-adapted transfer learning models such as RequiBERT [30]. Across these evaluations, DNA+ demonstrates consistent empirical superiority, delivering substantial improvements in classification accuracy and macro-averaged F1-scores while requiring orders of magnitude less computational overhead and zero dedicated GPU hardware during routine inference.

In summary, this study provides four primary scientific and technical contributions to automated software requirements engineering:
First, it introduces a dynamic syntactic-semantic dinucleotide codon grammar that models requirement statements as coupled primary-auxiliary dependencies, resolving the contextual limitations of earlier static single-nucleotide representations.
Second, it formulates an alignment-free multimodal feature manifold that synthesizes genuine biological genomic sequences, extracting twenty-three structural sequence features, information entropy, and dinucleotide relative abundance ratios, combined with lexical term frequencies and dense Sentence-BERT contextual embeddings.
Third, it delivers a rigorous empirical benchmark across three distinct requirements corpora comprising 8,654 authentic specifications, demonstrating statistically verified superiority, exceptional precision-recall retention under extreme class imbalance, and complete feature ablation transparency.
Fourth, it establishes a dual-tier explainable artificial intelligence framework uniting global Shapley feature importance with local token attribution, providing verified interpretability that bridges bio-inspired sequence processing and human requirements engineering practice.

The remainder of this article is structured as follows. Section 2 delineates the materials, benchmark corpus characteristics, and mathematical formulations of the dynamic dinucleotide codon grammar, sequence synthesis, and multimodal feature extraction pipeline. Section 3 presents the experimental setup, classification results, diagnostic evaluation curves, component ablation findings, and inferential statistical significance tests. Section 4 discusses the empirical findings in comparison with contemporary state-of-the-art literature, explores practical software engineering implications, addresses threats to validity, and outlines future research trajectories. Finally, Section 5 concludes the paper.